\section{Introduction}

\subsection{Motivation}
The position and orientation estimation has been a problem known in aviation over the years. With the advent of automated control systems, the need to determine exact location and orientation has steadily increased. This challenge extends beyond aviation to mechanics and multi-body systems, where precise positioning is essential for research, navigation, and control.\\

In mechanisms, particularly in multi-body mechanisms, accurate position and orientation measurements form the foundation for effective operation. However, not all systems are equipped with appropriate sensors, such as encoders, due to structural or economic limitations. This lack of direct measurement capabilities necessitates alternative methods for estimation, particularly in constrained environments.\\

Currently, there is no miniaturized and standalone measurement method capable of providing moderately good results in such systems. Traditional solutions either require external references, such as GPS or optical tracking, or rely on complex and costly setups unsuitable for many applications. The absence of a compact and self-sufficient solution highlights the need for an alternative approach.\\

Inertial sensors present a promising avenue for addressing these challenges. They are highly robust, capable of functioning in almost any environment, and do not depend on external infrastructure. However, their inherent inaccuracies necessitate advanced processing techniques to enhance their reliability. The integration of inertial sensors with a mechanism’s structural constraints offers a potential pathway to improving estimation accuracy in constrained multi-body systems.

\subsection{Literature review}

By nature, sensors have finite resolution and sampling time. Each sensor measurement is subject to errors, but many factors contribute to this, like bias or high-frequency noise. Some of these factors are inevitable, while others can be strongly reduced, e.g., mounting bias. To improve results, all sensors should be checked and calibrated before usage \cite{mi13060879} \cite{Hol_2011} \cite{gyro_calib}. It is also a good practice to pre-filter raw measurements with frequency filtration \cite{BADRI20101425} to minimize noise and characteristic disturbances. \\

Sensors have various measurement periods and modes. Especially inertial sensors, are incredibly fast, leaving only a couple of milliseconds for calculation. To achieve good-quality results, as stated in Nyquist–Shannon sampling theorem \cite{sample_theorem}, the frequency of estimation should be at least twice as high as the maximal frequency observed in the mechanism's movement. Achieving such high performance on an embedded system requires suitable implementation solutions. Since most of the calculations are conducted on floating precision numbers, hardware acceleration is required  \cite{fpu2} \cite{fpu}.\\

A huge impact on sensor fusion was caused by adopting the Kalman Filter in estimation. The mathematical concept of a filter based on the equation of state and statistics was first presented in 1960 \cite{kalman}. The result of widespread familiarity with the Kalman filter is also the development of many modifications and improvements that allow it to be used in specific applications. One of the most useful features presented in the thesis is the possibility of developing the filter with an additional correction to meet the constraints imposed on the system state variables. The concept and implementation of correction are detailed in the article \cite{simon}. The correction leads to an exact solution in the case of linear equalities or inequalities, whereas it gives an estimation in the case of non-linear constraints. It should be noted that, to the author's best knowledge, this method has not been utilized in estimating the position of robots and mechanisms. Although the method should probably be well adopted due to its universality, it needs to be verified.\\

The navigation methods outlined also appear in robotics, but the dominant part is mobile robotics \cite{accelerometer_mobile}. In stationary mechanisms, higher accuracy and repeatability are required. For this reason, a common approach is to use encoder readout and position estimation through a forward kinematics task \cite{forward_kinematics}. The significant advantage of this solution is that the received results comply with mechanism constraints. In the absence of an encoder or similar built-in sensor (e.g,. in biorobotics), computer vision algorithms \cite{cv_positioning} \cite{cv_positioning2} and triangulation methods \cite{igps} are used as a substitute. Inertial sensors are also mounted in industrial robots, but the measurements are usually used in quality and health checks like vibration measures \cite{Dogrusoz_2020}.\\

\subsection{Problem formulation}

The problem under consideration pertains to the development of a method for determining the position and orientation of a selected system component based on independent inertial sensors and additional system information. Traditional approaches to localization rely on direct position feedback, such as encoders or optical tracking systems, which are often impractical due to cost, complexity, or system constraints. This study seeks to provide a robust alternative using inertial sensors combined with system-specific constraints. \\

A fundamental challenge in this problem is that the additional system information cannot be directly applied through the selection of an alternative coordinate system. This limitation arises because the number of degrees of freedom may change over time. Consequently, traditional transformation techniques that rely on fixed constraints and pre-defined coordinate systems are unsuitable. Unlike conventional methods that predefine coordinate systems, this approach dynamically integrates constraints into the estimation process. This allows for adaptation to varying degrees of freedom without requiring fundamental changes to the underlying mathematical model. The proposed method utilizes an improved Extended Kalman Filter (EKF) with constraint correction, ensuring better stability and accuracy in position and orientation estimation compared to traditional inertial navigation approaches. By leveraging an iterative correction method inspired by the Newton-Raphson algorithm, the system can adjust to changes in constraints while maintaining robust state estimation.\\

The structure of the remaining content of this article is as follows. Section .... In Section , .... .... are discussed in Section . Finally, in Section \ref{sec:conclusions}, the conclusions are drawn.
